\documentclass{article}
\usepackage{graphicx}
\usepackage{amsmath}

\title{Quiz 3 - Laboratorio Integrado de Física}
\author{Juan Carlos Gómez Jaramillo \\c.c. 71217699}
\date{\today}

\begin{document}
	
	\maketitle
	
	\section{Preguntas}
	\begin{enumerate}
		\item ¿Explique en sus propias palabras el fenómeno de interferencia de ondas? 
		
		\vspace{0.3cm}
		
		Es un fenómeno que se produce cuando dos o más ondas coinciden en la misma región del espacio y de tiempo, sus perturbaciones se superponen dando como resultado, generando una onda  cuya amplitud depende de la combinación de la amplitud de las ondas originales, esta combinación puede aumentar o disminuir la amplitud total estas ultimas continúan su trayectoria de forma independiente tras la interacción. 
		
		\vspace{0.3cm}
		
		\item ¿Explique en sus propias palabras cómo varía la intensidad de una señal o de la radiación de una fuente puntual?
		
		\vspace{0.3cm}
		
		\vspace{0.3cm}
		
		\item ¿Cuáles son los objetivos que se deben cumplir en esta práctica?
		
		\vspace{0.3cm}
		
		\begin{itemize}
			\item Analizar el comportamiento de la señal WIFI del hogar, usando la aplicación Physics Toolbox Suite, para ello vamos a estudiar cómo varía la señal dentro de la casa haciendo mediciones con el sensor WIFI de la app, tomando datos durante 20 segundos en cada punto y anotando el máximo y el mínimo de intensidad de la señal, se va a reportar el promedio como valor de la medida como error: $\frac{\text{máx}-\text{mín}}{2}$..
			\item Comprender el funcionamiento del módem:
			Se pretende entender la forma en que el módem emite la señal y cómo esta disminuye su intensidad ante la presencia de distancia y obstáculos. Al posicionar el dispositivo a distancias controladas y conocidas, se evalúan las variaciones de la señal para verificar si los valores máximos y mínimos concuerdan con la teoría, incorporando para ello los principios de la ley del inverso del cuadrado de la distancia.
			\item Caracterizar la absorción de la señal WIFI por diferentes materiales:
			Se pretende evaluar cuánto se atenúa la señal según el material utilizado (ladrillo, vidrio, madera, metal y plástico). Mediante mediciones a una distancia fija entre el módem y el dispositivo, con y sin el material interpuesto, se aísla la caída de intensidad provocada por el obstáculo. Al replicar la prueba a otra distancia, se espera corroborar la consistencia de los coeficientes de atenuación calculados.
			\item Cuantificar la señal WIFI de todo el hogar:
			Se pretende mapear la distribución espacial de la señal a partir de los datos de una cuadrícula de medición, representando el promedio por punto con una escala de color e indicando la posición del módem. Teóricamente, un entorno sin obstáculos presentaría un patrón de círculos concéntricos con una intensidad proporcional a $1/r^2$. Las zonas que se apartan de este comportamiento permiten identificar los obstáculos del espacio. Finalmente, se contrasta la señal medida en los diferentes puntos del modelo teórico, comprobando que la diferencia sea coherente con los valores de atenuación medidos en el tercer objetivo.    
		\end{itemize}
		
		\vspace{0.3cm}
		
		\item Identifique al menos tres aplicaciones de las ondas electromagnéticas, que usted utiliza en su vida diaria.
		
		\vspace{0.3cm}
		
		\begin{itemize}
			\item Telecomunicaciones:
			Los teléfonos celulares y el módem, intercambian información mediante ondas electromagnéticas (wifi ~2,4 - 5 GHz, datos móviles ~ 700 MHz - 3,5 GHz). La información se transmite modulando la onda portadora. Su intensidad decae con la distancia según la ley del inverso cuadrado y se atenúa al atravesar paredes u otros obstáculos.
			
			\item Aplicaciones de salud:
			\textsl{Rayos X}$(10^{16} - 10^{19} \text{ Hz, alta energía ionizantes})$ 
			
			Radiografías: los huesos absorben más rayos X que los tejidos blandos, generando una imagen de contraste sobre una placa o sensor digital.
			
			\textsl{Tomografía computarizada} (TC o CT sacan): múltiples radiografías tomadas en ángulos distintos alrededor del cuerpo, reconstruidas por computador en imágenes 3D transversales. Da mucho más detalle que una radiografía simple.
			
			\textsl{Radioterapia}: rayos X de alta energía dirigidos para destruir células cancerígenas dañando su ADN.
			
			\item Desinfección: Lámparas UV-C para desinfectar instrumental médico, salas y agua, ya que dañan el ADN de bacterias y virus.
		\end{itemize}
		
		\vspace{0.3cm}
		
		\item Realice la propagación de incertidumbre apropiada para la función \[a=\frac{sin{(b\theta)}}{b\theta}\]
		
		\vspace{0.3cm}
		
		\vspace{0.3cm}
		
		
	\end{enumerate}
	
\end{document}